
\documentclass[aspectratio=169,11pt]{beamer}

\usetheme{AnnArbor}
\usepackage[utf8]{inputenc}
\usepackage[T1]{fontenc}
\usepackage{amsmath,amssymb,booktabs,array,multicol}
\usepackage{listings}
\usepackage{xcolor}
\usepackage{graphicx}
\usepackage{hyperref}

\definecolor{CodeBg}{RGB}{248,248,248}
\definecolor{TitleBlue}{RGB}{38,79,129}
\definecolor{AccentGreen}{RGB}{35,125,93}
\definecolor{SoftGray}{RGB}{80,80,80}
\definecolor{QuizOrange}{RGB}{180,95,30}

\setbeamercolor{title}{fg=TitleBlue}
\setbeamercolor{frametitle}{fg=TitleBlue}
\setbeamercolor{structure}{fg=AccentGreen}
\setbeamercolor{block title}{fg=white,bg=TitleBlue}
\setbeamercolor{block body}{bg=blue!3}
\setbeamercolor{alertblock title}{fg=white,bg=QuizOrange}
\setbeamercolor{alertblock body}{bg=orange!6}
\setbeamertemplate{navigation symbols}{}
\setbeamertemplate{footline}[frame number]
\setbeamertemplate{itemize item}{\raise0.15ex\hbox{\tiny$\blacktriangleright$}}
\setbeamertemplate{itemize subitem}{\raise0.15ex\hbox{\tiny$\bullet$}}

\lstdefinestyle{py}{
  language=Python,
  basicstyle=\ttfamily\scriptsize,
  backgroundcolor=\color{CodeBg},
  frame=single,
  framerule=0.2pt,
  rulecolor=\color{gray!35},
  breaklines=true,
  showstringspaces=false,
  keywordstyle=\color{TitleBlue}\bfseries,
  commentstyle=\color{SoftGray}\itshape,
  stringstyle=\color{AccentGreen}
}

\newcommand{\key}[1]{\textbf{\textcolor{AccentGreen}{#1}}}
\newcommand{\smallnote}[1]{\vspace{0.25em}{\footnotesize\color{SoftGray}#1}}
\newcommand{\quiztitle}[1]{\textbf{\textcolor{QuizOrange}{#1}}}
\newcommand{\imageplaceholder}[2]{%
\begin{center}
  \fcolorbox{SoftGray}{gray!5}{%
  \begin{minipage}[c][0.55\textheight][c]{0.84\linewidth}
  \centering
  {\Large\textcolor{TitleBlue}{Image Placeholder}}\\\[1.0em]
  {\normalsize\textbf{#2}}\\\[1.2em]
  {\footnotesize Source path in Markdown:}\\\[0.35em]
  {\small\texttt{\detokenize{#1}}}
  \end{minipage}}
\end{center}
}

\title[Introduction to NumPy]{Introduction to NumPy}
\subtitle{Arrays, Vectorization, Broadcasting, and Numerical Computing}
\author{Duc-Minh Vu}
\institute{SLSCM Lab -- FDA -- NEU}
\date{\today}

\begin{document}

\begin{frame}
  \titlepage
\end{frame}

\begin{frame}{Lesson Roadmap}
\begin{columns}[T]
\column{0.48\linewidth}
\begin{enumerate}
  \item What NumPy is
  \item Why NumPy is useful
  \item Installation and importing
  \item Creating arrays
  \item Array properties
  \item Indexing and slicing
  \item Reshaping, resizing, stacking, splitting
\end{enumerate}
\column{0.48\linewidth}
\begin{enumerate}
  \setcounter{enumi}{7}
  \item Broadcasting
  \item Arithmetic and aggregation
  \item Universal functions
  \item Boolean operations
  \item Linear algebra
  \item Random numbers and statistics
  \item Integration and common errors
\end{enumerate}
\end{columns}
\end{frame}

\begin{frame}{Learning Outcomes}
\small
By the end of this lesson, learners will be able to:
\begin{itemize}
  \item Explain the purpose of NumPy in numerical computing.
  \item Distinguish NumPy arrays from Python lists.
  \item Create one-dimensional and multidimensional arrays.
  \item Inspect \texttt{shape}, \texttt{size}, \texttt{ndim}, and \texttt{dtype}.
  \item Use indexing, slicing, reshaping, stacking, and splitting.
  \item Apply vectorized arithmetic, broadcasting, aggregation, and universal functions.
  \item Perform basic linear algebra and descriptive statistics.
  \item Generate random values and understand reproducibility.
  \item Explain how NumPy connects with Pandas, SciPy, and Scikit-learn.
\end{itemize}
\end{frame}

\begin{frame}[fragile]{Prerequisites and Setup}
\begin{block}{Recommended background}
\begin{itemize}
  \item Basic Python variables, lists, loops, and functions.
  \item A working Python environment: Jupyter Notebook, JupyterLab, Google Colab, or similar.
\end{itemize}
\end{block}
\begin{block}{Installation and import}
\begin{lstlisting}[style=py]
pip install numpy
\end{lstlisting}
\begin{lstlisting}[style=py]
import numpy as np
print(np.__version__)
\end{lstlisting}
\end{block}
\end{frame}


\begin{frame}{How to Read NumPy Commands}
\small
Most NumPy commands follow the same pattern:
\[
\texttt{np.function(input, option1=value1, option2=value2)}
\]
\begin{itemize}
  \item \texttt{np}: the standard alias for the NumPy package.
  \item \texttt{function}: the operation we want NumPy to perform.
  \item \texttt{input}: the array, list, shape, or numerical values passed to the function.
  \item \texttt{option=value}: named arguments that control the behavior of the command.
\end{itemize}
\begin{block}{Example}
\texttt{np.arange(0, 10, 2)} creates values starting at 0, stopping before 10, with step size 2.
\end{block}
\end{frame}

\begin{frame}{Command Pattern: Object, Attribute, Method, Function}
\small
\begin{tabular}{p{0.26\linewidth}p{0.62\linewidth}}
\toprule
\textbf{Form} & \textbf{Meaning} \\
\midrule
\texttt{np.array([1,2,3])} & Calls a NumPy function to create an array. \\
\texttt{arr.shape} & Reads an attribute that describes the array. \\
\texttt{arr.reshape(2,3)} & Calls a method attached to an existing array object. \\
\texttt{np.mean(arr)} & Calls a NumPy function that summarizes an array. \\
\texttt{arr.mean()} & Equivalent method-style version for many common operations. \\
\bottomrule
\end{tabular}
\vspace{0.5em}
\begin{alertblock}{Teaching note}
Students should identify three things in every command: input, operation, and output.
\end{alertblock}
\end{frame}

\begin{frame}[fragile]{Mini Example: Explain Each Line}
\small
\begin{lstlisting}[style=py]
import numpy as np

scores = np.array([7, 8, 6, 9])
bonus = scores + 1
average = scores.mean()
print(bonus)
print(average)
\end{lstlisting}
\begin{itemize}
  \item \texttt{import numpy as np}: loads NumPy and gives it the short name \texttt{np}.
  \item \texttt{np.array([...])}: converts a Python list into a NumPy array.
  \item \texttt{scores + 1}: adds 1 to every element using vectorization.
  \item \texttt{scores.mean()}: calculates the average value of the array.
  \item \texttt{print(...)}: displays the result.
\end{itemize}
\end{frame}

\section{What Is NumPy?}

\begin{frame}{What Is NumPy?}
\small
\key{NumPy}, short for \key{Numerical Python}, is a core Python library for fast numerical computing.
\vspace{0.4em}

Its central object is the \key{N-dimensional array}, called \texttt{ndarray}. A NumPy array can represent:
\begin{itemize}
  \item A one-dimensional vector.
  \item A two-dimensional matrix.
  \item A three-dimensional tensor.
  \item Higher-dimensional numerical structures.
\end{itemize}
\begin{block}{Why it matters}
NumPy stores homogeneous data efficiently and supports fast operations over whole arrays.
\end{block}
\end{frame}

\begin{frame}{NumPy Provides}
\small
\begin{columns}[T]
\column{0.47\linewidth}
\begin{itemize}
  \item Fast array operations.
  \item Vectorized calculations.
  \item Broadcasting.
  \item Linear algebra functions.
\end{itemize}
\column{0.47\linewidth}
\begin{itemize}
  \item Statistical functions.
  \item Random-number generation.
  \item Integration with Pandas, SciPy, Matplotlib, and Scikit-learn.
\end{itemize}
\end{columns}
\vspace{0.7em}
% \imageplaceholder{numpy\_linear\_algebra.webp}{NumPy linear algebra overview}
\end{frame}

\begin{frame}{Main Features}
\scriptsize
\begin{tabular}{p{0.26\linewidth}p{0.66\linewidth}}
\toprule
\textbf{Feature} & \textbf{Meaning} \\\\
\midrule
\texttt{ndarray} & Efficient multidimensional array with one common data type. \\\\
Vectorization & Apply calculations to entire arrays without explicit Python loops. \\\\
Broadcasting & Operate on compatible arrays with different shapes. \\\\
Linear algebra & Matrix multiplication, determinants, inverses, eigenvalues, and vector products. \\\\
Statistics & Mean, median, variance, standard deviation, percentiles, and quantiles. \\\\
Integration & Works with Pandas, SciPy, Matplotlib, Scikit-learn, and Statsmodels. \\\\
\bottomrule
\end{tabular}
\end{frame}

\begin{frame}{Quick Check: What Is NumPy?}
\small
\quiztitle{Question 1.} What is the main data structure in NumPy?
\begin{multicols}{2}
\begin{itemize}
  \item A. \texttt{DataFrame}
  \item B. \texttt{ndarray}
  \item C. \texttt{dictionary}
  \item D. \texttt{tuple}
\end{itemize}
\end{multicols}
\vspace{0.8em}
\quiztitle{Question 2. True or false?} NumPy is designed mainly for numerical computing.
\end{frame}

\section{Why Learn NumPy?}

\begin{frame}{Why Learn NumPy?}
\small
NumPy provides an efficient foundation for numerical work in Python.
\begin{itemize}
  \item Executes vectorized operations much faster than many ordinary Python loops.
  \item Stores homogeneous numerical data more compactly than standard Python lists.
  \item Provides optimized functions for matrices, linear algebra, and transformations.
  \item Supports random-number generation and statistical analysis.
  \item Expresses complex mathematical operations using concise syntax.
  \item Serves as the numerical foundation for many data-science libraries.
\end{itemize}
\end{frame}

\begin{frame}[fragile]{Python List vs. NumPy Array}
\begin{columns}[T]
\column{0.48\linewidth}
\textbf{Python list: mixed objects}
\begin{lstlisting}[style=py]
values = [10, 2.5, "Python", True]
\end{lstlisting}
\begin{itemize}
  \item Flexible object container.
  \item Not optimized for large numerical arrays.
\end{itemize}
\column{0.48\linewidth}
\textbf{NumPy array: homogeneous values}
\begin{lstlisting}[style=py]
import numpy as np
values = np.array([10, 20, 30, 40])
\end{lstlisting}
\begin{itemize}
  \item Regular structure.
  \item Efficient numerical computation.
\end{itemize}
\end{columns}
\end{frame}

\begin{frame}[fragile]{Example: Multiplying Values}
\begin{columns}[T]
\column{0.48\linewidth}
\textbf{Python list}
\begin{lstlisting}[style=py]
values = [1, 2, 3, 4]
result = []

for value in values:
    result.append(value * 10)

print(result)
# [10, 20, 30, 40]
\end{lstlisting}
\column{0.48\linewidth}
\textbf{NumPy vectorization}
\begin{lstlisting}[style=py]
import numpy as np
values = np.array([1, 2, 3, 4])

result = values * 10
print(result)
# [10 20 30 40]
\end{lstlisting}
\end{columns}
\begin{block}{Key idea}
The same operation is applied to all elements at once.
\end{block}
\end{frame}

\begin{frame}{Quick Check: Why Learn NumPy?}
\small
\quiztitle{Question 1.} Why are NumPy arrays efficient for numerical calculations?
\begin{itemize}
  \item A. They always contain text
  \item B. They use a homogeneous and optimized array structure
  \item C. They do not use memory
  \item D. They automatically connect to the internet
\end{itemize}
\vspace{0.6em}
\quiztitle{Question 2.} Which expression multiplies every element of a NumPy array \texttt{a} by 10?
\begin{multicols}{2}
\begin{itemize}
  \item A. \texttt{a * 10}
  \item B. \texttt{a.append(10)}
  \item C. \texttt{a.add("10")}
  \item D. \texttt{a.sort(10)}
\end{itemize}
\end{multicols}
\end{frame}

\section{Creating Arrays}

\begin{frame}[fragile]{Creating NumPy Arrays}
\small
NumPy arrays can be created from Python lists, tuples, ranges, or built-in NumPy functions.
\begin{lstlisting}[style=py]
import numpy as np

a = [9, 3, 3, 5]
arr = np.array(a)
print(arr)
# [9 3 3 5]
\end{lstlisting}
\vspace{0.2em}
\begin{columns}[T]
\column{0.5\linewidth}
\begin{lstlisting}[style=py]
arr = np.array([10, 20, 30, 40])
print(arr)
\end{lstlisting}
\column{0.5\linewidth}
\begin{lstlisting}[style=py]
matrix = np.array([
    [1, 2, 3],
    [4, 5, 6]
])
print(matrix)
\end{lstlisting}
\end{columns}
\end{frame}

\begin{frame}[fragile]{Multidimensional Arrays}
\small
A three-dimensional array can represent stacked matrices or tensors.
\begin{lstlisting}[style=py]
tensor = np.array([
    [
        [1, 2],
        [3, 4]
    ],
    [
        [5, 6],
        [7, 8]
    ]
])
print(tensor)
\end{lstlisting}
\begin{block}{Interpretation}
Dimensions often correspond to observations, features, time, channels, or scenarios.
\end{block}
\end{frame}

\begin{frame}[fragile]{Common Array-Creation Functions}
\small
\begin{columns}[T]
\column{0.48\linewidth}
\begin{lstlisting}[style=py]
np.zeros(5)
np.zeros((2, 3))
np.ones((2, 3))
np.full((2, 3), 7)
\end{lstlisting}
\column{0.48\linewidth}
\begin{lstlisting}[style=py]
np.arange(0, 10, 2)
# [0 2 4 6 8]

np.linspace(0, 1, 5)
# [0.   0.25 0.5  0.75 1.]

np.eye(3)
\end{lstlisting}
\end{columns}
\end{frame}

\begin{frame}{Key Commands: Creating Arrays}
\small
\begin{tabular}{p{0.28\linewidth}p{0.61\linewidth}}
\toprule
\textbf{Command} & \textbf{Meaning and when to use it} \\
\midrule
\texttt{np.array(data)} & Convert a Python list, tuple, or nested sequence into an \texttt{ndarray}. \\
\texttt{np.zeros(shape)} & Create an array of zeros; useful for initialization. \\
\texttt{np.ones(shape)} & Create an array of ones; useful for masks, weights, or initialization. \\
\texttt{np.full(shape, value)} & Create an array filled with one chosen value. \\
\texttt{np.arange(start, stop, step)} & Generate regularly spaced values; \texttt{stop} is excluded. \\
\texttt{np.linspace(start, stop, n)} & Generate exactly \texttt{n} evenly spaced values, including endpoints by default. \\
\texttt{np.eye(n)} & Create an \(n\times n\) identity matrix. \\
\bottomrule
\end{tabular}
\end{frame}

\begin{frame}{Quick Check: Creating Arrays}
\small
\quiztitle{Question 1.} Which function converts a Python list into a NumPy array?
\begin{multicols}{2}
\begin{itemize}
  \item A. \texttt{np.array()}
  \item B. \texttt{np.list()}
  \item C. \texttt{np.convert()}
  \item D. \texttt{np.ndarray\_list()}
\end{itemize}
\end{multicols}
\vspace{0.4em}
\quiztitle{Question 2.} Which function creates evenly spaced values between two endpoints?
\begin{multicols}{2}
\begin{itemize}
  \item A. \texttt{np.linspace()}
  \item B. \texttt{np.stack()}
  \item C. \texttt{np.split()}
  \item D. \texttt{np.mean()}
\end{itemize}
\end{multicols}
\end{frame}

\section{Array Properties}

\begin{frame}[fragile]{Array Properties: Inspecting an Array}
\small
\begin{lstlisting}[style=py]
arr = np.array([
    [1, 2, 3],
    [4, 5, 6]
])

print(arr.ndim)     # 2
print(arr.shape)    # (2, 3)
print(arr.size)     # 6
print(arr.dtype)
print(arr.itemsize)
\end{lstlisting}
\begin{block}{How to read these commands}
Each command asks NumPy to describe the array: number of dimensions, shape, total size, data type, and memory per element.
\end{block}
\end{frame}

\begin{frame}{Array Properties: Meaning of Each Attribute}
\small
\begin{tabular}{p{0.22\linewidth}p{0.68\linewidth}}
\toprule
\textbf{Attribute} & \textbf{Meaning} \\
\midrule
\texttt{ndim} & Number of dimensions. A vector has 1 dimension; a matrix has 2 dimensions. \\
\texttt{shape} & Size along each dimension. For example, \texttt{(2, 3)} means 2 rows and 3 columns. \\
\texttt{size} & Total number of elements in the array. \\
\texttt{dtype} & Data type stored in the array, such as integer or floating-point number. \\
\texttt{itemsize} & Number of bytes used by each element. \\
\bottomrule
\end{tabular}
\end{frame}

\begin{frame}{Key Commands: Inspecting Array Properties}
\small
\begin{tabular}{p{0.25\linewidth}p{0.64\linewidth}}
\toprule
\textbf{Attribute} & \textbf{What it tells us} \\
\midrule
\texttt{arr.ndim} & Number of axes/dimensions in the array. \\
\texttt{arr.shape} & Length of the array along each axis; e.g., \texttt{(3,4)} means 3 rows and 4 columns. \\
\texttt{arr.size} & Total number of stored elements. \\
\texttt{arr.dtype} & Type used to store each element, such as \texttt{int64} or \texttt{float64}. \\
\texttt{arr.itemsize} & Number of bytes used by one element. \\
\texttt{arr.nbytes} & Total bytes used by all array elements. \\
\bottomrule
\end{tabular}
\begin{block}{Useful habit}
Before reshaping, broadcasting, or modeling, inspect \texttt{shape} and \texttt{dtype}; many NumPy errors originate from unexpected shapes or types.
\end{block}
\end{frame}

\begin{frame}{Quick Check: Array Properties}
\small
\quiztitle{Question 1.} Which attribute returns the dimensions of an array?
\begin{multicols}{2}
\begin{itemize}
  \item A. \texttt{shape}
  \item B. \texttt{mean}
  \item C. \texttt{append}
  \item D. \texttt{index}
\end{itemize}
\end{multicols}
\vspace{0.4em}
\quiztitle{Question 2.} Which attribute returns the total number of elements?
\begin{multicols}{2}
\begin{itemize}
  \item A. \texttt{size}
  \item B. \texttt{dtype}
  \item C. \texttt{ndim}
  \item D. \texttt{itemsize}
\end{itemize}
\end{multicols}
\end{frame}

\section{Indexing and Slicing}

\begin{frame}[fragile]{Array Indexing}
\begin{columns}[T]
\column{0.48\linewidth}
\textbf{One-dimensional indexing}
\begin{lstlisting}[style=py]
arr = np.array([10, 20, 30, 40])

print(arr[0])   # 10
print(arr[2])   # 30
print(arr[-1])  # 40
\end{lstlisting}
\column{0.48\linewidth}
\textbf{Two-dimensional indexing}
\begin{lstlisting}[style=py]
matrix = np.array([
    [1, 2, 3],
    [4, 5, 6]
])

print(matrix[0, 1])  # 2
print(matrix[1, 2])  # 6
\end{lstlisting}
\end{columns}
\begin{block}{Indexing convention}
Python and NumPy use zero-based indexing.
\end{block}
\end{frame}

\begin{frame}[fragile]{Array Slicing}
\begin{columns}[T]
\column{0.48\linewidth}
\textbf{One-dimensional slicing}
\begin{lstlisting}[style=py]
arr = np.array([10, 20, 30, 40, 50])

print(arr[1:4])
# [20 30 40]

print(arr[::2])
# [10 30 50]
\end{lstlisting}
\column{0.48\linewidth}
\textbf{Two-dimensional slicing}
\begin{lstlisting}[style=py]
matrix = np.array([
    [1, 2, 3],
    [4, 5, 6],
    [7, 8, 9]
])

print(matrix[0:2, 1:3])
# [[2 3]
#  [5 6]]
\end{lstlisting}
\end{columns}
\end{frame}

\begin{frame}[fragile]{Views and Copies}
\small
Many NumPy slices return a \key{view} of the original array instead of an independent copy.
\begin{lstlisting}[style=py]
arr = np.array([10, 20, 30, 40])

view = arr[1:3]
view[0] = 999

print(arr)
# [ 10 999  30  40]
\end{lstlisting}
Create an independent copy when needed:
\begin{lstlisting}[style=py]
copy = arr[1:3].copy()
\end{lstlisting}
\begin{alertblock}{Important}
Modifying a view may modify the original array.
\end{alertblock}
\end{frame}

\begin{frame}{Key Commands: Indexing and Slicing}
\small
\begin{tabular}{p{0.28\linewidth}p{0.61\linewidth}}
\toprule
\textbf{Syntax} & \textbf{Meaning} \\
\midrule
\texttt{arr[i]} & Select the element at position \texttt{i}; negative indices count from the end. \\
\texttt{arr[a:b]} & Select positions \texttt{a} through \texttt{b-1}. \\
\texttt{arr[a:b:s]} & Slice with step size \texttt{s}. \\
\texttt{matrix[i,j]} & Select row \texttt{i}, column \texttt{j}. \\
\texttt{matrix[r1:r2,c1:c2]} & Extract a rectangular subarray. \\
\texttt{arr[mask]} & Boolean indexing: retain values where \texttt{mask} is \texttt{True}. \\
\texttt{arr.copy()} & Make independent data when a view should not modify the original. \\
\bottomrule
\end{tabular}
\end{frame}

\begin{frame}{Quick Check: Indexing and Slicing}
\small
\quiztitle{Question 1.} What does \texttt{arr[0]} return?
\begin{multicols}{2}
\begin{itemize}
  \item A. The first element
  \item B. The last element
  \item C. The array shape
  \item D. The array size
\end{itemize}
\end{multicols}
\quiztitle{Question 2.} In a two-dimensional array, \texttt{matrix[1, 2]} refers to:
\begin{itemize}
  \item A. Row index 1 and column index 2
  \item B. Row 1 only
  \item C. Column 2 only
  \item D. The array dimensions
\end{itemize}
\quiztitle{Question 3. True or false?} A NumPy slice may share memory with the original array.
\end{frame}

\section{Reshaping, Stacking, Splitting}

\begin{frame}[fragile]{Reshaping Arrays}
\small
Reshaping changes array dimensions without changing its data.
\begin{lstlisting}[style=py]
arr = np.arange(12)
matrix = arr.reshape(3, 4)
print(matrix)
# [[ 0  1  2  3]
#  [ 4  5  6  7]
#  [ 8  9 10 11]]
\end{lstlisting}
Use \texttt{-1} to infer a dimension:
\begin{lstlisting}[style=py]
matrix = arr.reshape(2, -1)
\end{lstlisting}
Flatten an array:
\begin{lstlisting}[style=py]
flat = matrix.flatten()
flat_view = matrix.ravel()
\end{lstlisting}
\smallnote{\texttt{flatten()} returns a copy; \texttt{ravel()} often returns a view when possible.}
\end{frame}

\begin{frame}[fragile]{Resizing Arrays}
\small
\texttt{resize()} changes shape and may change the number of elements.
\begin{lstlisting}[style=py]
arr = np.array([1, 2, 3, 4])
resized = np.resize(arr, (2, 3))
print(resized)
# [[1 2 3]
#  [4 1 2]]
\end{lstlisting}
\begin{block}{Compare}
\begin{itemize}
  \item \texttt{reshape()} requires the total number of elements to remain compatible.
  \item \texttt{resize()} can create a different total size by repeating values.
\end{itemize}
\end{block}
\end{frame}

\begin{frame}[fragile]{Stacking Arrays}
\begin{columns}[T]
\column{0.5\linewidth}
\begin{lstlisting}[style=py]
a = np.array([1, 2, 3])
b = np.array([4, 5, 6])

np.vstack((a, b))
# [[1 2 3]
#  [4 5 6]]
\end{lstlisting}
\column{0.5\linewidth}
\begin{lstlisting}[style=py]
np.hstack((a, b))
# [1 2 3 4 5 6]

np.stack((a, b), axis=0)
\end{lstlisting}
\end{columns}
\begin{block}{Purpose}
Stacking combines arrays into larger arrays along a chosen axis.
\end{block}
\end{frame}

\begin{frame}[fragile]{Splitting Arrays}
\small
Splitting divides arrays into smaller arrays.
\begin{lstlisting}[style=py]
arr = np.array([1, 2, 3, 4, 5, 6])
parts = np.split(arr, 3)
print(parts)
# [array([1, 2]), array([3, 4]), array([5, 6])]
\end{lstlisting}
For multidimensional arrays:
\begin{lstlisting}[style=py]
matrix = np.array([
    [1, 2, 3, 4],
    [5, 6, 7, 8]
])

left, right = np.hsplit(matrix, 2)
\end{lstlisting}
\end{frame}

\begin{frame}{Key Commands: Reshape, Stack, and Split}
\scriptsize
\begin{tabular}{p{0.29\linewidth}p{0.60\linewidth}}
\toprule
\textbf{Command} & \textbf{Meaning} \\
\midrule
\texttt{arr.reshape(r,c)} & Reorganize the same elements into a new compatible shape. \\
\texttt{arr.reshape(...,-1)} & Ask NumPy to infer one dimension automatically. \\
\texttt{arr.flatten()} & Return a flattened 1-D \emph{copy}. \\
\texttt{arr.ravel()} & Return a flattened view when possible. \\
\texttt{np.resize(arr, shape)} & Change total size if needed; values may repeat. \\
\texttt{np.vstack((a,b))} & Stack arrays vertically (row direction). \\
\texttt{np.hstack((a,b))} & Stack arrays horizontally. \\
\texttt{np.stack((a,b), axis=k)} & Add a new axis and stack along that axis. \\
\texttt{np.split(arr,n)} & Divide an array into \texttt{n} equal parts when possible. \\
\texttt{np.hsplit(matrix,n)} & Split a matrix along the column direction. \\
\bottomrule
\end{tabular}
\end{frame}

\begin{frame}{Quick Check: Reshaping and Splitting}
\small
\quiztitle{Question 1.} Which method changes an array from one shape to another?
\begin{multicols}{2}
\begin{itemize}
  \item A. \texttt{reshape()}
  \item B. \texttt{mean()}
  \item C. \texttt{split()}
  \item D. \texttt{sort\_index()}
\end{itemize}
\end{multicols}
\quiztitle{Question 2.} What does \texttt{-1} mean in \texttt{reshape()}?
\begin{itemize}
  \item A. NumPy should infer that dimension
  \item B. Delete one dimension
  \item C. Reverse the array
  \item D. Convert values to negative numbers
\end{itemize}
\quiztitle{Question 3.} Which function stacks arrays vertically? \texttt{np.vstack()}, \texttt{np.mean()}, \texttt{np.split()}, or \texttt{np.random()}?
\end{frame}

\section{Broadcasting and Arithmetic}

\begin{frame}[fragile]{Broadcasting: Core Idea}
\small
Broadcasting allows NumPy to perform operations between arrays of compatible shapes.
\begin{lstlisting}[style=py]
arr = np.array([1, 2, 3, 4])
result = arr + 10
print(result)
# [11 12 13 14]
\end{lstlisting}
\begin{block}{Basic rule}
Starting from the rightmost dimensions, two dimensions are compatible when:
\begin{itemize}
  \item they are equal; or
  \item one of them is \texttt{1}.
\end{itemize}
\end{block}
\end{frame}

\begin{frame}[fragile]{Broadcasting Between Arrays}
\small
\begin{lstlisting}[style=py]
matrix = np.array([
    [1, 2, 3],
    [4, 5, 6]
])

row = np.array([10, 20, 30])
result = matrix + row
print(result)
# [[11 22 33]
#  [14 25 36]]
\end{lstlisting}
\begin{alertblock}{Incompatible shapes}
\begin{lstlisting}[style=py]
a = np.ones((2, 3))
b = np.ones((2, 2))
# a + b raises a broadcasting error
\end{lstlisting}
\end{alertblock}
\end{frame}

\begin{frame}[fragile]{Basic Arithmetic Operations}
\small
NumPy supports element-wise arithmetic.
\begin{lstlisting}[style=py]
a = np.array([1, 2, 3])
b = np.array([4, 5, 6])

print(a + b)   # [5 7 9]
print(a - b)   # [-3 -3 -3]
print(a * b)   # [ 4 10 18]
print(a / b)
print(a ** 2)  # [1 4 9]
print(b % a)
\end{lstlisting}
\begin{block}{Important distinction}
\texttt{\*} performs element-wise multiplication, not matrix multiplication.
\end{block}
\end{frame}

\begin{frame}{Key Commands: Arithmetic and Broadcasting}
\small
\begin{tabular}{p{0.28\linewidth}p{0.61\linewidth}}
\toprule
\textbf{Expression} & \textbf{Meaning} \\
\midrule
\texttt{a + b}, \texttt{a - b} & Element-wise addition and subtraction. \\
\texttt{a * b} & Element-wise multiplication. \\
\texttt{a / b} & Element-wise division. \\
\texttt{a ** p} & Raise every element to power \texttt{p}. \\
\texttt{a \% b} & Element-wise remainder. \\
\texttt{array + scalar} & Broadcast a scalar to every element. \\
\texttt{matrix + row} & Broadcast a compatible 1-D row across matrix rows. \\
\bottomrule
\end{tabular}
\begin{alertblock}{Check shapes first}
Broadcasting compares dimensions from right to left. Dimensions must be equal or one of them must be 1.
\end{alertblock}
\end{frame}

\begin{frame}{Quick Check: Broadcasting and Arithmetic}
\small
\quiztitle{Question 1.} What does broadcasting allow?
\begin{itemize}
  \item A. Operations on arrays with compatible different shapes
  \item B. Automatic internet transmission
  \item C. Conversion of arrays into files
  \item D. Removal of all dimensions
\end{itemize}
\quiztitle{Question 2.} Two dimensions are broadcasting-compatible when they are equal or:
\begin{multicols}{2}
\begin{itemize}
  \item A. One of them is 1
  \item B. Both are negative
  \item C. Both are text
  \item D. Their sum is zero
\end{itemize}
\end{multicols}
\quiztitle{Question 3.} What does \texttt{a * b} do for arrays of the same shape?
\end{frame}

\section{Aggregation and Universal Functions}

\begin{frame}[fragile]{Aggregation Functions}
\small
Aggregation functions summarize an array.
\begin{lstlisting}[style=py]
arr = np.array([9, 3, 3, 5])

print(arr.sum())   # 20
print(arr.mean())  # 5.0
print(arr.min())
print(arr.max())
print(np.median(arr))
print(np.var(arr))
print(np.std(arr))
\end{lstlisting}
\begin{block}{Typical use}
Aggregation turns raw arrays into summary statistics.
\end{block}
\end{frame}

\begin{frame}[fragile]{Aggregation Along an Axis}
\small
\begin{lstlisting}[style=py]
matrix = np.array([
    [1, 2, 3],
    [4, 5, 6]
])

print(matrix.sum(axis=0))
# [5 7 9]

print(matrix.sum(axis=1))
# [ 6 15]
\end{lstlisting}
\begin{itemize}
  \item \texttt{axis=0}: aggregate down rows, producing column-wise results.
  \item \texttt{axis=1}: aggregate across columns, producing row-wise results.
\end{itemize}
\end{frame}

\begin{frame}[fragile]{Universal Functions}
\small
A universal function, or \key{ufunc}, performs element-wise operations on arrays.
\begin{columns}[T]
\column{0.48\linewidth}
\begin{lstlisting}[style=py]
arr = np.array([1, 4, 9, 16])
print(np.sqrt(arr))
# [1. 2. 3. 4.]

print(np.exp(np.array([0, 1, 2])))
\end{lstlisting}
\column{0.48\linewidth}
\begin{lstlisting}[style=py]
values = np.array([1, np.e, np.e**2])
print(np.log(values))

angles = np.array([0, np.pi / 2, np.pi])
print(np.sin(angles))
\end{lstlisting}
\end{columns}
\vspace{0.4em}
Other examples: \texttt{np.abs()}, \texttt{np.round()}.
\end{frame}

\begin{frame}{Key Commands: Aggregation and Universal Functions}
\scriptsize
\begin{tabular}{p{0.29\linewidth}p{0.60\linewidth}}
\toprule
\textbf{Command} & \textbf{Meaning} \\
\midrule
\texttt{arr.sum()}, \texttt{arr.mean()} & Sum and arithmetic mean. \\
\texttt{arr.min()}, \texttt{arr.max()} & Minimum and maximum. \\
\texttt{np.median(arr)} & Median. \\
\texttt{np.var(arr)}, \texttt{np.std(arr)} & Variance and standard deviation. \\
\texttt{func(..., axis=0)} & Aggregate down rows, producing one result per column. \\
\texttt{func(..., axis=1)} & Aggregate across columns, producing one result per row. \\
\texttt{np.sqrt(arr)} & Square root element by element. \\
\texttt{np.exp(arr)}, \texttt{np.log(arr)} & Exponential and natural logarithm element by element. \\
\texttt{np.sin(arr)} & Sine element by element; angles are in radians. \\
\texttt{np.abs(arr)}, \texttt{np.round(arr)} & Absolute value and rounding. \\
\bottomrule
\end{tabular}
\end{frame}

\begin{frame}{Quick Check: Aggregation and Ufuncs}
\small
\quiztitle{Question 1.} Which function calculates the arithmetic mean?
\begin{multicols}{2}
\begin{itemize}
  \item A. \texttt{np.mean()}
  \item B. \texttt{np.stack()}
  \item C. \texttt{np.dot()}
  \item D. \texttt{np.reshape()}
\end{itemize}
\end{multicols}
\quiztitle{Question 2.} In a two-dimensional array, what does \texttt{axis=0} generally aggregate over?
\begin{itemize}
  \item A. Rows, producing column-wise results
  \item B. Columns, producing row-wise results
  \item C. File names
  \item D. Data types
\end{itemize}
\quiztitle{Question 3.} Which function calculates square roots element by element?
\end{frame}

\section{Boolean Operations}

\begin{frame}[fragile]{Comparison and Boolean Filtering}
\small
NumPy supports element-wise comparisons.
\begin{lstlisting}[style=py]
arr = np.array([10, 20, 30, 40])
print(arr > 20)
# [False False  True  True]

selected = arr[arr > 20]
print(selected)
# [30 40]
\end{lstlisting}
Combine conditions:
\begin{lstlisting}[style=py]
selected = arr[(arr >= 20) & (arr <= 30)]
print(selected)
# [20 30]
\end{lstlisting}
Use \texttt{\&} for AND, \texttt{|} for OR, and \texttt{\textasciitilde} for NOT.
\end{frame}

\begin{frame}{Key Commands: Boolean Operations}
\small
\begin{tabular}{p{0.31\linewidth}p{0.58\linewidth}}
\toprule
\textbf{Expression} & \textbf{Meaning} \\
\midrule
\texttt{arr > value} & Produce a Boolean array from an element-wise comparison. \\
\texttt{arr[arr > value]} & Filter and return only matching values. \\
\texttt{(cond1) \& (cond2)} & Element-wise logical AND. \\
\texttt{(cond1) | (cond2)} & Element-wise logical OR. \\
\texttt{\textasciitilde mask} & Element-wise logical NOT. \\
\texttt{np.where(cond,a,b)} & Choose \texttt{a} where condition is true and \texttt{b} otherwise. \\
\bottomrule
\end{tabular}
\begin{block}{Important syntax}
Use parentheses around combined comparisons because \texttt{\&} and \texttt{|} have different precedence from comparison operators.
\end{block}
\end{frame}

\begin{frame}{Quick Check: Boolean Operations}
\small
\quiztitle{Question 1.} What does \texttt{arr[arr > 20]} return?
\begin{itemize}
  \item A. Elements greater than 20
  \item B. The array size
  \item C. The array data type
  \item D. All elements converted to Boolean values
\end{itemize}
\vspace{0.8em}
\quiztitle{Discussion.} Why is Boolean filtering useful in data analysis?
\end{frame}

\section{Linear Algebra}

\begin{frame}[fragile]{Matrix Multiplication}
\small
\begin{lstlisting}[style=py]
A = np.array([
    [1, 2],
    [3, 4]
])

result = np.dot(A, A)
print(result)
# [[ 7 10]
#  [15 22]]

result = A @ A
\end{lstlisting}
\begin{block}{Element-wise vs. matrix multiplication}
\begin{itemize}
  \item \texttt{A \* A}: element-wise multiplication.
  \item \texttt{A @ A}: matrix multiplication.
\end{itemize}
\end{block}
\end{frame}

\begin{frame}[fragile]{Matrix Operations}
\small
\begin{columns}[T]
\column{0.48\linewidth}
\begin{lstlisting}[style=py]
print(A.T)

determinant = np.linalg.det(A)
print(determinant)

inverse = np.linalg.inv(A)
print(inverse)
\end{lstlisting}
\column{0.48\linewidth}
\begin{lstlisting}[style=py]
A = np.array([
    [2, 1],
    [1, 3]
])

b = np.array([8, 13])
x = np.linalg.solve(A, b)
print(x)
\end{lstlisting}
\end{columns}
\begin{block}{Condition}
An inverse exists only for a square, nonsingular matrix.
\end{block}
\end{frame}

\begin{frame}[fragile]{Eigenvalues and Vector Products}
\small
\begin{columns}[T]
\column{0.48\linewidth}
\begin{lstlisting}[style=py]
eigenvalues, eigenvectors = np.linalg.eig(A)

print(eigenvalues)
print(eigenvectors)
\end{lstlisting}
\column{0.48\linewidth}
\begin{lstlisting}[style=py]
a = np.array([1, 2, 3])
b = np.array([4, 5, 6])

print(np.inner(a, b))
print(np.outer(a, b))
print(np.dot(a, b))
print(np.vdot(a, b))
\end{lstlisting}
\end{columns}
\smallnote{For real one-dimensional arrays, \texttt{dot} and \texttt{vdot} are often the same. \texttt{vdot} also handles complex conjugation.}
\end{frame}

\begin{frame}{Key Commands: Linear Algebra}
\scriptsize
\begin{tabular}{p{0.29\linewidth}p{0.60\linewidth}}
\toprule
\textbf{Command} & \textbf{Meaning} \\
\midrule
\texttt{A @ B} & Matrix multiplication; preferred readable syntax. \\
\texttt{np.dot(A,B)} & Dot product / matrix product depending on dimensions. \\
\texttt{A.T} & Transpose an array. \\
\texttt{np.linalg.det(A)} & Compute determinant of a square matrix. \\
\texttt{np.linalg.inv(A)} & Compute inverse of a nonsingular square matrix. \\
\texttt{np.linalg.solve(A,b)} & Solve \(Ax=b\) directly; usually preferable to computing \(A^{-1}b\). \\
\texttt{np.linalg.eig(A)} & Return eigenvalues and eigenvectors. \\
\texttt{np.inner(a,b)} & Inner product. \\
\texttt{np.outer(a,b)} & Outer product. \\
\texttt{np.vdot(a,b)} & Vector dot product with complex conjugation when needed. \\
\bottomrule
\end{tabular}
\end{frame}

\begin{frame}{Quick Check: Linear Algebra}
\small
\quiztitle{Question 1.} Which operator performs matrix multiplication?
\begin{multicols}{2}
\begin{itemize}
  \item A. \texttt{@}
  \item B. \texttt{\%}
  \item C. \texttt{//}
  \item D. \texttt{\&}
\end{itemize}
\end{multicols}
\quiztitle{Question 2.} Which function calculates a matrix inverse?
\begin{itemize}
  \item A. \texttt{np.linalg.inv()}
  \item B. \texttt{np.mean()}
  \item C. \texttt{np.resize()}
  \item D. \texttt{np.random()}
\end{itemize}
\quiztitle{Question 3.} Explain the difference between \texttt{A * B} and \texttt{A @ B}.
\end{frame}

\section{Random Numbers and Statistics}

\begin{frame}[fragile]{Random-Number Generation}
\small
For new code, NumPy commonly uses a random-number generator object.
\begin{lstlisting}[style=py]
rng = np.random.default_rng()
\end{lstlisting}
\begin{columns}[T]
\column{0.48\linewidth}
\begin{lstlisting}[style=py]
rng = np.random.default_rng(42)
values = rng.integers(
    low=1,
    high=10,
    size=5
)
print(values)
\end{lstlisting}
\column{0.48\linewidth}
\begin{lstlisting}[style=py]
values = rng.normal(
    loc=0,
    scale=1,
    size=5
)
print(values)
\end{lstlisting}
\end{columns}
\begin{alertblock}{Note}
NumPy random-number generators are for numerical and statistical work; they should not be treated as cryptographically secure generators.
\end{alertblock}
\end{frame}

\begin{frame}[fragile]{Common Distributions}
\small
\begin{lstlisting}[style=py]
rng = np.random.default_rng(42)

rng.uniform(low=0, high=1, size=5)
rng.normal(loc=0, scale=1, size=5)
rng.binomial(n=10, p=0.5, size=5)
rng.poisson(lam=3, size=5)
rng.exponential(scale=2, size=5)
rng.chisquare(df=4, size=5)
\end{lstlisting}
\begin{block}{Reproducibility}
Using the same seed helps reproduce the same pseudorandom sequence in a compatible environment.
\end{block}
\end{frame}

\begin{frame}[fragile]{Statistical Functions}
\small
\begin{lstlisting}[style=py]
data = np.array([12, 15, 18, 20, 25])

print(np.mean(data))
print(np.median(data))
print(np.var(data))
print(np.std(data))
print(np.percentile(data, 25))
print(np.quantile(data, 0.75))
\end{lstlisting}
\begin{block}{Population and sample conventions}
By default, \texttt{np.var()} and \texttt{np.std()} use \texttt{ddof=0}. For the common sample convention, use \texttt{ddof=1}.
\end{block}
\end{frame}

\begin{frame}[fragile]{Missing Numerical Values}
\small
NumPy commonly represents missing numerical values with \texttt{np.nan}.
\begin{lstlisting}[style=py]
data = np.array([10.0, 20.0, np.nan, 40.0])
print(np.mean(data))
# nan

print(np.nanmean(data))
print(np.nanmedian(data))
print(np.nanstd(data))

print(np.isnan(data))
clean_data = data[~np.isnan(data)]
\end{lstlisting}
\smallnote{More complete missing-data workflows are commonly handled with Pandas.}
\end{frame}

\begin{frame}{Key Commands: Random Numbers and Statistics}
\scriptsize
\begin{tabular}{p{0.32\linewidth}p{0.57\linewidth}}
\toprule
\textbf{Command} & \textbf{Meaning} \\
\midrule
\texttt{np.random.default\_rng(seed)} & Create a modern random-number generator; a seed supports reproducibility. \\
\texttt{rng.integers(low,high,size)} & Generate random integers; \texttt{high} is excluded. \\
\texttt{rng.uniform(...)} & Sample from a uniform distribution. \\
\texttt{rng.normal(loc,scale,size)} & Sample from a normal distribution with mean \texttt{loc} and SD \texttt{scale}. \\
\texttt{rng.binomial(...)} / \texttt{poisson(...)} & Sample from common discrete distributions. \\
\texttt{np.percentile(data,q)} & Return percentile \texttt{q} on a 0--100 scale. \\
\texttt{np.quantile(data,q)} & Return quantile \texttt{q} on a 0--1 scale. \\
\texttt{np.std(data,ddof=1)} & Sample-standard-deviation convention. \\
\texttt{np.isnan(data)} & Identify \texttt{NaN} values. \\
\texttt{np.nanmean(data)} & Mean while ignoring \texttt{NaN}. \\
\bottomrule
\end{tabular}
\end{frame}

\begin{frame}{Quick Check: Random Numbers and Statistics}
\small
\quiztitle{Question 1.} Which method creates a modern NumPy random-number generator?
\begin{itemize}
  \item A. \texttt{np.random.default\_rng()}
  \item B. \texttt{np.random.file()}
  \item C. \texttt{np.create\_random\_array()}
  \item D. \texttt{np.random.text()}
\end{itemize}
\quiztitle{Question 2.} Why is a seed useful?
\begin{itemize}
  \item A. It helps reproduce a pseudorandom sequence
  \item B. It makes values truly unpredictable
  \item C. It removes every random value
  \item D. It converts data into strings
\end{itemize}
\quiztitle{Question 3.} Which parameter gives the usual sample standard deviation convention? \texttt{ddof=1}, \texttt{axis=-100}, \texttt{dtype="sample"}, or \texttt{copy=False}?
\end{frame}

\section{Vectorization and Performance}

\begin{frame}[fragile]{Vectorized Operations}
\small
Vectorization applies an operation to an entire array at once.
\begin{columns}[T]
\column{0.48\linewidth}
\begin{lstlisting}[style=py]
a = np.arange(5)
result = a * 10
print(result)
# [ 0 10 20 30 40]
\end{lstlisting}
\column{0.48\linewidth}
\begin{lstlisting}[style=py]
a = list(range(5))
result = []
for value in a:
    result.append(value * 10)
print(result)
\end{lstlisting}
\end{columns}
\begin{block}{Typical benefits}
Shorter code, easier reading, faster operations for large numerical arrays, and better fit for scientific workflows.
\end{block}
\end{frame}

\begin{frame}[fragile]{Vectorized Conditional Logic}
\small
\begin{lstlisting}[style=py]
arr = np.array([10, 20, 30, 40])

labels = np.where(
    arr >= 30,
    "high",
    "low"
)

print(labels)
# ['low' 'low' 'high' 'high']
\end{lstlisting}
\begin{block}{Interpretation}
Vectorized conditions are useful for labeling, screening, filtering, and transforming datasets.
\end{block}
\end{frame}

\begin{frame}{Key Commands: Vectorization}
\small
\begin{tabular}{p{0.30\linewidth}p{0.59\linewidth}}
\toprule
\textbf{Pattern} & \textbf{Meaning} \\
\midrule
\texttt{arr * scalar} & Apply the arithmetic operation to every element without an explicit Python loop. \\
\texttt{np.sqrt(arr)} & Apply a ufunc element by element in compiled NumPy code. \\
\texttt{arr >= threshold} & Evaluate a condition for the whole array at once. \\
\texttt{np.where(cond,a,b)} & Vectorized conditional transformation. \\
\bottomrule
\end{tabular}
\begin{block}{Core idea}
Vectorization describes \emph{how we express computation}: operations are written for arrays rather than manually iterating over individual elements.
\end{block}
\end{frame}

\begin{frame}{Quick Check: Vectorization}
\small
\quiztitle{Question 1.} What is a major advantage of vectorization?
\begin{itemize}
  \item A. It applies array operations without explicit Python loops
  \item B. It converts every array into text
  \item C. It eliminates the need for memory
  \item D. It prevents all errors
\end{itemize}
\vspace{0.8em}
\quiztitle{Discussion.} When might vectorized code be easier to maintain than loop-based code?
\end{frame}

\section{Memory, Sorting, Images}

\begin{frame}[fragile]{Memory Considerations}
\small
NumPy arrays usually store homogeneous values in a regular memory layout.
\begin{lstlisting}[style=py]
arr = np.array([1, 2, 3, 4], dtype=np.int32)
print(arr.nbytes)

small_values = np.array([1, 2, 3, 4], dtype=np.int8)
print(small_values.dtype)
print(small_values.nbytes)
\end{lstlisting}
\begin{block}{Data type choice}
A smaller data type may reduce memory use, but it must represent the required value range and precision.
\end{block}
\end{frame}

\begin{frame}[fragile]{Converting Data Types}
\small
\begin{lstlisting}[style=py]
arr = np.array([1.2, 2.8, 3.5])
integers = arr.astype(np.int32)
print(integers)
# [1 2 3]
\end{lstlisting}
\begin{alertblock}{Important}
Conversion from floating-point values to integers removes the fractional part.
\end{alertblock}
\end{frame}

\begin{frame}[fragile]{Sorting and Searching}
\small
\begin{lstlisting}[style=py]
arr = np.array([9, 3, 7, 1])
sorted_arr = np.sort(arr)
print(sorted_arr)
# [1 3 7 9]

arr = np.array([10, 20, 30, 40])
indices = np.where(arr > 20)
print(indices)

arr = np.array([1, 2, 2, 3, 3, 3])
values, counts = np.unique(arr, return_counts=True)
print(values)
print(counts)
\end{lstlisting}
\end{frame}

\begin{frame}[fragile]{Working with Images}
\small
Digital images can be represented as NumPy arrays.
\begin{itemize}
  \item Grayscale image: two-dimensional array.
  \item Color image: three-dimensional array containing height, width, and color channels.
\end{itemize}
\begin{lstlisting}[style=py]
image = np.array([
    [0, 128, 255],
    [255, 128, 0],
    [50, 100, 150]
])
print(image.shape)
# (3, 3)

brighter = np.clip(image + 30, 0, 255)
\end{lstlisting}
\smallnote{\texttt{np.clip()} keeps values inside a valid range.}
\end{frame}

\begin{frame}{Key Commands: Memory, Types, Sorting, and Search}
\scriptsize
\begin{tabular}{p{0.30\linewidth}p{0.59\linewidth}}
\toprule
\textbf{Command} & \textbf{Meaning} \\
\midrule
\texttt{arr.nbytes} & Total memory used by the array elements. \\
\texttt{arr.astype(dtype)} & Convert array values to a new data type. \\
\texttt{np.sort(arr)} & Return sorted values. \\
\texttt{np.where(condition)} & Return indices satisfying a condition when only the condition is supplied. \\
\texttt{np.unique(arr)} & Return unique values. \\
\texttt{np.unique(arr, return\_counts=True)} & Return unique values and their frequencies. \\
\texttt{np.clip(arr,low,high)} & Limit every value to a specified range; useful for image intensities. \\
\bottomrule
\end{tabular}
\end{frame}

\begin{frame}{Quick Check: Memory, Search, Images}
\small
\quiztitle{Question 1.} Which attribute returns the number of bytes used by array elements?
\begin{multicols}{2}
\begin{itemize}
  \item A. \texttt{nbytes}
  \item B. \texttt{ndim}
  \item C. \texttt{mean}
  \item D. \texttt{shape}
\end{itemize}
\end{multicols}
\quiztitle{Question 2.} Which function returns unique array values?
\begin{itemize}
  \item A. \texttt{np.unique()}
  \item B. \texttt{np.reshape()}
  \item C. \texttt{np.random()}
  \item D. \texttt{np.outer()}
\end{itemize}
\quiztitle{Question 3.} A grayscale image is commonly represented as what type of array?
\end{frame}

\section{Integration and Errors}

\begin{frame}[fragile]{Integration with Pandas}
\small
A Pandas \texttt{Series} or \texttt{DataFrame} can exchange data with NumPy.
\begin{lstlisting}[style=py]
import numpy as np
import pandas as pd

arr = np.array([
    [1, 2],
    [3, 4],
    [5, 6]
])

df = pd.DataFrame(arr, columns=["A", "B"])
array_from_df = df.to_numpy()

df["A_sqrt"] = np.sqrt(df["A"])
\end{lstlisting}
\end{frame}

\begin{frame}[fragile]{Integration with Scikit-learn}
\small
Many Scikit-learn models accept NumPy arrays as input.
\begin{lstlisting}[style=py]
from sklearn.linear_model import LinearRegression

X = np.array([[1], [2], [3], [4]])
y = np.array([2, 4, 6, 8])

model = LinearRegression()
model.fit(X, y)

prediction = model.predict(np.array([[5]]))
print(prediction)
\end{lstlisting}
NumPy supports feature matrices, target vectors, model predictions, numerical preprocessing, and evaluation calculations.
\end{frame}

\begin{frame}[fragile]{Common NumPy Errors}
\small
\begin{columns}[T]
\column{0.48\linewidth}
\textbf{Shape mismatch}
\begin{lstlisting}[style=py]
a = np.ones((2, 3))
b = np.ones((2, 2))
# a + b raises an error
\end{lstlisting}
\textbf{Invalid reshape}
\begin{lstlisting}[style=py]
arr = np.arange(10)
# arr.reshape(3, 4) raises an error
\end{lstlisting}
\column{0.48\linewidth}
\textbf{Division by zero}
\begin{lstlisting}[style=py]
arr = np.array([1.0, 0.0])
print(1 / arr)
\end{lstlisting}
\textbf{Integer overflow}
\begin{lstlisting}[style=py]
arr = np.array([127], dtype=np.int8)
print(arr + 1)
\end{lstlisting}
\end{columns}
\end{frame}

\begin{frame}{Good Practices}
\small
\begin{itemize}
  \item Import NumPy using \texttt{import numpy as np}.
  \item Use vectorized operations instead of Python loops when practical.
  \item Check \texttt{shape}, \texttt{ndim}, and \texttt{dtype} before complex operations.
  \item Use broadcasting only when intended shape behavior is clear.
  \item Distinguish element-wise multiplication from matrix multiplication.
  \item Use \texttt{.copy()} when an independent array is required.
  \item Choose data types carefully to balance range, precision, and memory.
  \item Use a fixed random seed when reproducibility is important.
  \item Treat \texttt{NaN}, infinity, and overflow explicitly.
\end{itemize}
\end{frame}

\begin{frame}{Key Commands: Integration with Data-Science Tools}
\small
\begin{tabular}{p{0.32\linewidth}p{0.57\linewidth}}
\toprule
\textbf{Command} & \textbf{Role} \\
\midrule
\texttt{pd.DataFrame(arr,...)} & Build a Pandas DataFrame from a NumPy array. \\
\texttt{df.to\_numpy()} & Extract DataFrame values as a NumPy array. \\
\texttt{np.sqrt(df["A"])} & Apply a NumPy ufunc to compatible Pandas data. \\
\texttt{model.fit(X,y)} & Train a Scikit-learn model using array-like feature and target data. \\
\texttt{model.predict(Xnew)} & Produce predictions for new array-like observations. \\
\bottomrule
\end{tabular}
\begin{block}{Connection}
NumPy arrays are a common numerical interchange format across Pandas, SciPy, Matplotlib, and Scikit-learn.
\end{block}
\end{frame}

\begin{frame}{Quick Check: Integration and Errors}
\small
\quiztitle{Question 1.} Which Pandas method converts a DataFrame to a NumPy array?
\begin{multicols}{2}
\begin{itemize}
  \item A. \texttt{to\_numpy()}
  \item B. \texttt{to\_list\_only()}
  \item C. \texttt{as\_matrix\_text()}
  \item D. \texttt{convert\_numpy\_file()}
\end{itemize}
\end{multicols}
\quiztitle{Question 2.} In machine learning, a two-dimensional NumPy array commonly represents:
\begin{itemize}
  \item A. A feature matrix
  \item B. A filename
  \item C. A chart title
  \item D. A package installer
\end{itemize}
\quiztitle{Question 3.} Why does \texttt{np.arange(10).reshape(3, 4)} fail?
\end{frame}

\section{Summary and Practice}

\begin{frame}{Content Summary I}
\small
\begin{tabular}{p{0.25\linewidth}p{0.66\linewidth}}
\toprule
\textbf{Topic} & \textbf{Main idea} \\
\midrule
NumPy & Core Python library for numerical computing. \\
\texttt{ndarray} & N-dimensional homogeneous array. \\
Vectorization & Array operations without explicit Python loops. \\
Broadcasting & Operations on compatible different shapes. \\
Indexing and slicing & Access and extract array elements. \\
Reshaping & Change array dimensions. \\
\bottomrule
\end{tabular}
\end{frame}

\begin{frame}{Content Summary II}
\small
\begin{tabular}{p{0.25\linewidth}p{0.66\linewidth}}
\toprule
\textbf{Topic} & \textbf{Main idea} \\
\midrule
Aggregation & Sums, means, minima, maxima, variance, and standard deviation. \\
Universal functions & Fast element-wise mathematical functions. \\
Linear algebra & Matrix multiplication, inverse, determinant, and eigenvalues. \\
Random generation & Generate values from probability distributions. \\
Statistics & Mean, median, variance, standard deviation, and percentiles. \\
Integration & Works closely with Pandas, SciPy, and Scikit-learn. \\
\bottomrule
\end{tabular}
\end{frame}

\begin{frame}{Practice: Array Creation}
\small
\begin{block}{Exercise 1}
\begin{enumerate}
  \item Create an array containing values from 1 to 20.
  \item Reshape it into a \(4 \times 5\) matrix.
  \item Print its shape, number of dimensions, size, and data type.
\end{enumerate}
\end{block}
\begin{block}{Exercise 2}
Using a \(4 \times 5\) matrix:
\begin{enumerate}
  \item Extract the first row.
  \item Extract the last column.
  \item Extract the central \(2 \times 3\) section.
  \item Select all even values using Boolean filtering.
\end{enumerate}
\end{block}
\end{frame}

\begin{frame}{Practice: Statistics and Broadcasting}
\small
\begin{block}{Exercise 3: Statistics}
Create an array of ten numerical values and calculate mean, median, variance, sample standard deviation, minimum, maximum, and the 25th, 50th, and 75th percentiles.
\end{block}
\begin{block}{Exercise 4: Broadcasting}
\begin{enumerate}
  \item Create a \(3 \times 4\) matrix.
  \item Create a one-dimensional array with four values.
  \item Add the one-dimensional array to each row of the matrix.
  \item Explain why broadcasting is valid.
\end{enumerate}
\end{block}
\end{frame}

\begin{frame}{Practice: Linear Algebra and Missing Values}
\small
\begin{block}{Exercise 5: Linear algebra}
Given
\[
A=\begin{bmatrix}2&1\\ 1&3\end{bmatrix},\quad b=\begin{bmatrix}8\\ 13\end{bmatrix},
\]
calculate the determinant of \(A\), inverse of \(A\), solve \(Ax=b\), and verify with \texttt{A @ x}.
\end{block}
\begin{block}{Exercise 6: Missing values}
Given \texttt{[10.0, np.nan, 20.0, 30.0, np.nan, 40.0]}, count missing values, calculate the mean while ignoring missing values, remove missing values, and replace missing values with the median.
\end{block}
\end{frame}

\begin{frame}{Final Reflection}
\small
\begin{itemize}
  \item Which NumPy topic is most important for data analysis workflows?
  \item When should you check \texttt{shape} before running an operation?
  \item How do vectorization and broadcasting reduce manual loops?
  \item Why is it important to distinguish views from copies?
  \item When should you use \texttt{ddof=1} instead of the default?
  \item What common NumPy error are you most likely to encounter in practice?
\end{itemize}
\end{frame}

\end{document}